% Secciones 1--3 del delta gravitatorio íntegro.
\section{Dominio genealógico de la realización gravitatoria}
\label{sec:l9s102:gravedad-dominio-genealogico}

La realización gravitatoria recibe como objeto de partida una sección del
continuo discreto ya construida por la composición
\begin{equation}
 \mathrm{APP}\longrightarrow\mathrm{TRIT}\longrightarrow\mathrm{TPK}
 \longrightarrow X_{\mathrm{TPK}}^{\mathrm{enr}}
 \longrightarrow\mathcal C_{\mathrm{cont}}^{\mathrm{disc}}.
 \label{eq:l9s102:gravedad-cadena-genealogica}
\end{equation}
APP conserva las evaluaciones aditiva y multiplicativa mediante residuo,
cociente y acarreo; TRIT determina el régimen y la orientación; el TPK
incorpora el selector trítico de fase sin eliminar ninguna de las dos hojas
APP, efectúa el transporte cardinal, conserva la memoria y realiza el retorno
de fase sin reiniciar el estado. El transductor
dimensional actúa después sobre la sección superviviente del continuo:
\begin{equation}
 \mathfrak G_{108}:\mathcal C_{\mathrm{cont}}^{\mathrm{disc}}
 \times\mathcal L_S^+
 \longrightarrow\mathcal G_{\mathrm{HMT\text{-}MD}},
 \qquad
 (x,\hbar_a)\longmapsto
 \bigl(G_a,R_{108},m_{108,a},\ell_{P,a}^{,2}\bigr).
 \label{eq:l9s102:gravedad-mapa-realizacion}
\end{equation}
Los cardinales \(90/120\), la periodicidad \(108\), la acción y el par
angular entran con su genealogía conservada. Ninguna identificación
gravitatoria posterior selecciona retroactivamente estos antecedentes.

\section{Transductor gravitatorio y selector cronológico de periodicidad}
\label{sec:l9s102:gravedad-transductor}

Sean \(c_{\mathrm{int}}>0\), \(t_0>0\),
\(\ell_0=c_{\mathrm{int}}t_0\) y
\(\hbar_a\in\mathcal L_S^+\), con
\(a\in\{\mathrm{pre},\mathrm{ret}\}\). El transductor dimensional es
\begin{equation}
 G_a(\theta)
 =\theta\frac{c_{\mathrm{int}}^5t_0^2}{\hbar_a}
 =\theta\frac{c_{\mathrm{int}}^3\ell_0^2}{\hbar_a},
 \qquad \theta>0.
 \label{eq:l9s102:g-transductor}
\end{equation}
La cronologia ciclica de \(108\) estados determina
\begin{equation}
 T_{108}=108t_0,
 \qquad
 \omega_{108}=\frac{2\pi_{\rm HMT}}{108t_0},
 \qquad
 R_{108}=\frac{c_{\mathrm{int}}}{\omega_{108}}
 =\frac{54}{\pi_{\rm HMT}}\ell_0.
 \label{eq:l9s102:g-reloj108}
\end{equation}
El cuanto cronológico y su masa asociada son
\begin{equation}
 E_{108,a}=\hbar_a\omega_{108},
 \qquad
 m_{108,a}=\frac{E_{108,a}}{c_{\mathrm{int}}^2}.
 \label{eq:l9s102:g-cuanto-masa}
\end{equation}
La longitud de Compton reducida satisface
\begin{equation}
 \bar\lambda_{C,a}
 =\frac{\hbar_a}{m_{108,a}c_{\mathrm{int}}}=R_{108}.
 \label{eq:l9s102:g-compton}
\end{equation}

\begin{theorem}[Unicidad del selector gravitatorio de periodicidad]
\label{thm:l9s102:g-selector}
Con la convencion reducida \(r_g=Gm/c_{\mathrm{int}}^2\), la condición
\(r_{g,a}(\theta)=R_{108}\) posee la solución positiva única
\begin{equation}
 \boxed{
 \theta_{108}^{\mathrm{circ}}=\left(\frac{54}{\pi_{\rm HMT}}\right)^2.}
 \label{eq:l9s102:g-theta}
\end{equation}
La realización correspondiente identifica estructuralmente
\begin{equation}
 R_{108}=\bar\lambda_{C,a}=r_{g,a}=\ell_P^{\mathrm{circ}}.
 \label{eq:l9s102:g-cuatro-longitudes}
\end{equation}
\end{theorem}

\begin{proof}
Por \cref{eq:l9s102:g-transductor,eq:l9s102:g-cuanto-masa},
\[
 r_{g,a}(\theta)
 =\frac{G_a(\theta)m_{108,a}}{c_{\mathrm{int}}^2}
 =\theta\frac{\pi_{\rm HMT}}{54}\ell_0.
\]
La igualdad con \(R_{108}=(54/\pi_{\rm HMT})\ell_0\) impone
\(\theta=(54/\pi_{\rm HMT})^2\). La linealidad en \(\theta\) y la positividad de
los factores determinan la unicidad.
\end{proof}

En la carta en la que \(c_{\mathrm{int}}=(\mu\varepsilon)^{-1/2}\), el
transductor adopta la forma
\begin{equation}
 \boxed{
 G_a=\left(\frac{54}{\pi_{\rm HMT}}\right)^2
 \frac{t_0^2}{\hbar_a(\mu\varepsilon)^{5/2}}.}
 \label{eq:l9s102:g-mu-epsilon}
\end{equation}
El producto constitutivo \(\mu\varepsilon\) gobierna el cono causal y la
escala gravitatoria; el cociente \(Z^2=\mu/\varepsilon\) conserva la
orientación de impedancia.

\section{Covariancia bidireccional de gravedad y acción}
\label{sec:l9s102:g-hbar}

Definimos
\begin{equation}
 R_{\mathrm{act}}=\frac{\hbar_{\mathrm{pre}}}
 {\hbar_{\mathrm{ret}}}.
 \label{eq:l9s102:r-act}
\end{equation}
La masa cronológica y el acoplamiento gravitatorio transforman
contragredientemente:
\begin{equation}
 \frac{m_{108,\mathrm{pre}}}{m_{108,\mathrm{ret}}}=R_{\mathrm{act}},
 \qquad
 \frac{G_{\mathrm{pre}}}{G_{\mathrm{ret}}}=R_{\mathrm{act}}^{-1}.
 \label{eq:l9s102:g-hbar-reciprocidad}
\end{equation}
Su producto permanece fijo:
\begin{equation}
 \boxed{
 G_a\hbar_a
 =\theta_{108}^{\mathrm{circ}}c_{\mathrm{int}}^5t_0^2,
 \qquad
 \ell_P^2=\frac{G_a\hbar_a}{c_{\mathrm{int}}^3}
 =\theta_{108}^{\mathrm{circ}}\ell_0^2.}
 \label{eq:l9s102:g-hbar-invariante}
\end{equation}
La acción conserva la orientación de la hoja; el área de Planck conserva la
geometría común de ambas orientaciones.

\begin{theorem}[Independencia de carta de las secciones \(h\) y \(G\)]
\label{thm:l9s102:covariancia-unidades-g}
Sean \(u_L,u_T,u_S\) bases positivas de las rectas dimensionales de
longitud, tiempo y acción, y cámbiense por
\[
 u_L'=\lambda_Lu_L,\qquad
 u_T'=\lambda_Tu_T,\qquad
 u_S'=\lambda_Su_S,
 \qquad \lambda_L,\lambda_T,\lambda_S>0.
\]
Las coordenadas de las mismas secciones verifican
\begin{equation}
 c_{\rm int}'=\frac{\lambda_T}{\lambda_L}c_{\rm int},\qquad
 t_0'=\frac{t_0}{\lambda_T},\qquad
 \hbar_a'=\frac{\hbar_a}{\lambda_S},\qquad
 G_a'=\frac{\lambda_S\lambda_T^3}{\lambda_L^5}G_a.
 \label{eq:l9s102:cambio-carta-g}
\end{equation}
Para la base inducida
\(u_G=u_L^5u_T^{-3}u_S^{-1}\) se obtiene
\begin{equation}
 \hbar_a'u_S'=\hbar_au_S,
 \qquad G_a'u_G'=G_au_G.
 \label{eq:l9s102:secciones-h-g-invariantes}
\end{equation}
Por consiguiente,
\(h_a=2\pi_{\rm HMT}\hbar_a\),
\(G_a=\theta_{108}^{\rm circ}c_{\rm int}^5t_0^2/\hbar_a\) y
\(G_a\hbar_a=\theta_{108}^{\rm circ}c_{\rm int}^5t_0^2\)
son identidades de secciones, no igualdades dependientes de una elección de
unidades. La carta metrológica publica sus coordenadas numéricas después de
la construcción HMT--MD.
\end{theorem}

\begin{proof}
Las tres primeras transformaciones de
\eqref{eq:l9s102:cambio-carta-g} son la ley contragrediente de coordenadas
en \(L\otimes T^{-1}\), \(T\) y \(S\). Su sustitución en
\eqref{eq:l9s102:g-transductor} produce la cuarta. La base de la recta
gravitatoria se transforma por el factor inverso, lo que prueba
\eqref{eq:l9s102:secciones-h-g-invariantes}. Los escalares
\(\pi_{\rm HMT}\) y \(\theta_{108}^{\rm circ}\) son adimensionales y
permanecen invariantes.
\end{proof}
